\documentclass[11pt,a4paper]{article}
\usepackage[margin=23mm]{geometry}
\usepackage{amsmath,amssymb,booktabs,microtype}
\usepackage[hidelinks]{hyperref}
\newcommand{\id}{\mathbb I}
\newcommand{\tr}{\operatorname{tr}}
\newcommand{\ket}[1]{|#1\rangle}
\newcommand{\bra}[1]{\langle#1|}
\newcommand{\pr}[1]{\ket{#1}\bra{#1}}
\title{Route G-R3: A Density Candidate and Two Relational Clocks}
\author{An explicit interaction, not a fitted slowdown law}
\date{24 September 2026}
\begin{document}
\maketitle
\begin{abstract}
We define an excitation-energy density relative to a declared counting
measure, couple it to one clock, and compute its readings relative to
another in an exact finite stationary state. The resulting relative
rate is observable under fixed calibration, but its sign is determined
by the interaction, not by the word density. A separate dispersive
exchange calculation admits both signs. Unobserved source fluctuations
produce a distribution of conditional phases, not generally a single
rate evaluated at the mean density. This closes one bounded candidate
test, not Chen's density--time claim, emergent gravity, or a universal
time-dilation principle.
\end{abstract}

\section{Define the object before assigning it dynamics}
Retain the G-R1 warning: energy statistics alone do not determine a
relational history. A candidate full information object would include
the joint state, event algebra, recording instruments and reference
measure. We do not identify any one scalar with all of that information.

Take an addressed matter subsystem $D$ with ground state $\ket0$,
excited state $\ket1$, and positive excitation energy $\epsilon_D$.
Fix the measure of this one operational cell to $\nu(D)=1$. Define
\begin{equation}\label{eq:density}
 N_D=\pr1,\qquad H_D-E_{D,g}\id=\epsilon_DN_D,\qquad
 \widehat E_d=\frac{\epsilon_DN_D}{\nu(D)},\qquad
 E_0=\frac{\epsilon_D}{\nu(D)} .
\end{equation}
Units are energy per declared cell, \emph{not} energy per spatial volume.
For several declared cells, an additive operator measure can be
$\mu_E(S)=\sum_{j\in S}\epsilon_jN_j$ with counting measure
$\nu(S)=|S|$; its average is $\mu_E(S)/\nu(S)$. The addressed set,
weights and coupling must be fixed together, not changed to fit a rate.
This one-cell test cannot empirically distinguish density from total
excitation energy. Ground subtraction prevents an arbitrary source
energy zero from changing the interaction.

Let $B$ be a three-level clock, $H_B=\epsilon N_B$ with
$N_B=\operatorname{diag}(0,1,2)$. Choose the single interaction
\begin{equation}\label{eq:interaction}
 V=\kappa H_B\frac{\widehat E_d}{E_0}
   =\kappa\epsilon N_BN_D,\qquad
 H_{BD}=\epsilon N_B(1+\kappa N_D)+\epsilon_DN_D .
\end{equation}
Tensor identities are suppressed. We impose $\kappa>-1$ for positive
clock gaps. $[N_D,H_{BD}]=0$: this is a dispersive, excitation-preserving
coupling. $\kappa$ is a declared coupling, not a fitted time function.
A reference clock $A$ is shielded from $D$. This asymmetry is an input.

\newpage
\section{One common constraint and an exact finite solution}
We explicitly extend the G-R1 stationary model to $A\otimes B\otimes D$.
We do not equate it with G-R2's recording-history constraint.
For an exactly reproducible card fix
\begin{equation}
 \epsilon_D=\epsilon,\quad \delta=\epsilon/4,\quad
 \kappa=1/4,\quad D_A=15,\quad E_*=14\delta,\quad
 H_A=\sum_{j=0}^{14}(14-j)\delta\,\pr j_A .
\end{equation}
The alternative audit cards $\kappa=0,-1/4$ use the \emph{same}
$A$ spectrum, calibration and readout. These rational choices make the
finite recurrence exact; they are not measurements or fits.
For $n=0,1,2$ and $d=0,1$, define
\begin{equation}
 e_{nd}=\epsilon n(1+\kappa d)+\epsilon_Dd,\qquad
 j_{nd}=e_{nd}/\delta=4n(1+\kappa d)+4d .
\end{equation}
All $j_{nd}$ are integers between $0$ and $14$ on these cards.
For normalized $a_n,c_d$, put
\begin{align}
 K&=H_A+H_{BD}-E_*\id,\\
 \ket\Psi&=\sum_{n,d}a_nc_d\ket{j_{nd}}_A\ket n_B\ket d_D,
 &K\ket\Psi&=0 .\label{eq:stationary}
\end{align}
Each term has total energy $E_*$; coincident $j_{nd}$ cause no
normalization problem because the $BD$ labels remain orthogonal.
Both local energies and the combined $H_{BD}$ are nonnegative.
The matched state is an assumed preparation, not a dynamically
derived attractor. No arbitrary incommensurate-spectrum extension is claimed.

Use the reference phase POVM
\begin{equation}
 \ket\phi_A=D_A^{-1/2}\sum_{j=0}^{14}e^{ij\phi}\ket j,\qquad
 F_A(d\phi)=\frac{D_A}{2\pi}\pr\phi_A\,d\phi .
\end{equation}
Its integral is $\id_A$; phase density is uniformly $1/(2\pi)$.
Normalized conditioning of \eqref{eq:stationary} gives
\begin{equation}\label{eq:conditional}
 \ket{\psi_{BD}(\phi)}
 =\sum_{n,d}a_nc_de^{-ie_{nd}\phi/\delta}\ket{n,d},
 \qquad i\delta\,\partial_\phi\ket{\psi_{BD}}=H_{BD}\ket{\psi_{BD}} .
\end{equation}
This comes directly from the constraint, without inserting an external
physical time. Only after calibration write $\tau_A=\hbar\phi/\delta$.
For a generic interaction involving the reference itself the simple
conditional equation need not hold \cite{smith}. Here $A$ is decoupled
from $BD$ except through the stationary correlation, so it holds exactly.

For completeness, if an event effect $F$ fails to commute with $K$,
the energy-sector pinching
\begin{equation}
 {\cal P}_K(F)=\sum_\ell P_\ell F P_\ell,\qquad
 K=\sum_\ell k_\ell P_\ell
\end{equation}
is positive, preserves POVM completeness, commutes with $K$, and gives
the same expectation in $\Psi$. The joint probabilities thus admit
constraint-invariant effects. This is not a microscopic implementation
of a nonlocal measurement apparatus.

\newpage
\section{Compare two readings, not two coordinate labels}
Set $a_n=1/\sqrt3$. Clock $B$ has phase states and POVM
\begin{equation}
 \ket\theta_B=\frac1{\sqrt3}\sum_{n=0}^2e^{-in\theta}\ket n,\qquad
 F_B(d\theta)=\frac3{2\pi}\pr\theta_B\,d\theta .
\end{equation}
The Fourier sign is opposite to that of $A$, whose Hamiltonian has the
complementary spectrum. Neither phase is a geometric rotation.
Conditional on source outcome $d$, \eqref{eq:conditional} becomes,
up to a source-only phase, $\ket{\theta_d(\phi)}_B$ with
\begin{equation}
 \theta_d(\phi)=\frac{\epsilon}{\delta}(1+\kappa d)\phi
 \pmod{2\pi}.
\end{equation}
The actual likelihood, rather than a zero-width clock reading, is
\begin{equation}\label{eq:likelihood}
 p(\theta\mid\phi,d)
 =\frac{1}{6\pi}\left|\sum_{n=0}^2
      e^{in[\theta-\theta_d(\phi)]}\right|^2
 =\frac{[1+2\cos(\theta-\theta_d)]^2}{6\pi}.
\end{equation}
It integrates to one. The complete joint density is
$p(\phi,\theta,d)=|c_d|^2p(\theta\mid\phi,d)/(2\pi)$.
On a fixed three-outcome grid $\theta_j=2\pi j/3$, the corresponding
projective probabilities are $[1+2\cos(\theta_j-\theta_d)]^2/9$.

Calibrate the uncoupled $B$ dial by $\tau_B=\hbar\theta/\epsilon$,
keeping that calibration when $D$ is excited. On a specified unwrapped
phase branch, the shift of the likelihood peak gives
\begin{equation}\label{eq:rate}
 r_{B/A}(d)=\frac{d\tau_B}{d\tau_A}
  =\frac{\delta}{\epsilon}\frac{d\theta_d}{d\phi}
  =1+\kappa d
  =1+\kappa\frac{E_d(d)}{E_0}.
\end{equation}
This is a result of \eqref{eq:interaction}, not an independent
phenomenological slowdown fit. At $d=1$ the three audit cards give
$r=3/4,1,5/4$: either direction is possible.
Finite POVM width requires repeated calibrated comparisons.
One event does not reveal an exact rate or a unique cycle count.
Use a predeclared short branch, or the event/epoch distinction of G-R2,
to avoid unwrapping ambiguity.

Two invariance checks matter. A common energy rescaling
$(\epsilon,\epsilon_D,\delta,E_*)\mapsto s(\epsilon,\epsilon_D,\delta,E_*)$
rescales $K$ but leaves $\Psi$, all phase probabilities and $r$ unchanged.
A passive relabeling of a dial, with its POVM and calibration transformed
consistently, also cannot change these probabilities. Recalibrating $B$
separately in every density sector would instead remove the comparison
one set out to measure.

For the same commuting conditional-generator structure, if both clocks
receive factors $1+\kappa_A d$ and $1+\kappa_B d$, then
\begin{equation}
 r_{B/A}=\frac{1+\kappa_Bd}{1+\kappa_Ad}.
\end{equation}
Equal factors cancel. This elementary common-mode diagnostic is not a
claim to have solved an arbitrary interaction with the finite reference.
Only a differential response is observable by the two-clock comparison.

\newpage
\section{Fluctuating density: why the mean is not one effective clock}
Let $w_d=|c_d|^2$. When $D$ is not read, the reduced clock state is
\begin{align}
 \rho_B(\phi)
  &=\sum_d w_d\pr{\theta_d(\phi)},\\
 [\rho_B(\phi)]_{nm}
  &=\frac13 e^{-i(n-m)x}
       \sum_d w_d e^{-i(n-m)\kappa d x},
 &x&=\epsilon\phi/\delta .\label{eq:characteristic}
\end{align}
The last factor is a characteristic function of the density distribution.
It is not generally $\exp[-i(n-m)\kappa\langle d\rangle x]$.
For small $\kappa x$,
\begin{equation}
 \log\langle e^{-iyd}\rangle
 =-iy\langle d\rangle-\tfrac12y^2\operatorname{Var}(d)+O(y^3),
 \qquad y=(n-m)\kappa x .
\end{equation}
The mean shifts phase, while the variance first changes visibility.
Replacing the operator by its mean drops precisely that information.
This respects, rather than overturns, the G-R1 density-only warning.

For $w_0=w_1=1/2$, the adjacent coherence is
\begin{equation}
 [\rho_B]_{10}=\tfrac13e^{-ix}
    \tfrac12(1+e^{-i\kappa x}).
\end{equation}
With $\kappa=1/4$ at $\phi=\pi$, $x=4\pi$; adjacent coherences vanish,
but $[\rho_B]_{20}=1/3$ and $\tr\rho_B^2=5/9$.
This is neither a pure clock with an average rate nor a fully dephased
state. A coherent pure source entangles with $B$; a classical source
mixture yields the same reduced $B$ statistics. These data alone do
not witness quantum source coherence.

An exact probability test uses $\phi=\pi/2$ on the same card.
Then $\theta_0=2\pi$ and $\theta_1=5\pi/2$. In the fixed $B$ grid:
\begin{center}
\begin{tabular}{lccc}
\toprule
source & $p_0$ & $p_1$ & $p_2$\\
\midrule
$d=0$ & $1$ & $0$ & $0$\\
$d=1$ & $1/9$ & $(1+\sqrt3)^2/9$ & $(1-\sqrt3)^2/9$\\
unread equal weights & $5/9$ & $(1+\sqrt3)^2/18$ & $(1-\sqrt3)^2/18$\\
\bottomrule
\end{tabular}
\end{center}
In contrast, replacing $d$ by $1/2$ in the Hamiltonian predicts
\begin{equation}
 p_0^{\rm mean\ substitution}
 =\frac{(1+\sqrt2)^2}{9}
 =\frac{3+2\sqrt2}{9}\neq\frac59 .
\end{equation}
This tests the approximation without fitting any data. The source
projector measurement itself commutes with $H_{BD}$, although
conditioning on its result selects a branch. A phase measurement on
the clock is a separate instrument with backreaction; it cannot be
silently replaced by nondisturbing, continuously available readings.

\newpage
\section{A microscopic sign check, separate from the exact rational card}
Use the Jaynes--Cummings exchange model, in frequency units
($\hbar=1$ in this section):
\begin{equation}
 H_{\rm JC}=\omega b^\dagger b+\Omega N_D
       +g(b\sigma_++b^\dagger\sigma_-),\qquad \Delta=\Omega-\omega .
\end{equation}
This is an established idealized model \cite{blais}, not a hardware
calibration. For $g\sqrt{n+1}/|\Delta|\ll1$, put
$X=b\sigma_+$ and $S=(g/\Delta)(X-X^\dagger)$. Direct algebra gives
\begin{equation}
 [S,H_0]=-V,\qquad
 [X,X^\dagger]=(2N_D-1)n+N_D,\qquad n=b^\dagger b .
\end{equation}
Hence the second-order transformed Hamiltonian is
\begin{align}
 H_{\rm eff}
 &=H_0+\tfrac12[S,V]+O(g^3/\Delta^2)\nonumber\\
 &=(\omega-\chi)n+(\Omega+\chi)N_D+2\chi nN_D
     +O(g^3/\Delta^2),\qquad \chi=\frac{g^2}{\Delta}.
\end{align}
Energy eigenvalue errors start at fourth order. With ground-sector
calibration $\omega_g=\omega-\chi>0$, the excited-sector rate is
\begin{equation}
 r_{e/g}^{(2)}=\frac{\omega+\chi}{\omega-\chi}
 =1+\kappa_{\rm eff},\qquad
 \kappa_{\rm eff}=\frac{2g^2/\Delta}{\omega-g^2/\Delta}.
\end{equation}
Positive and negative detuning yield opposite signs. The interaction
does not force denser matter to slow the clock.

The exact fixed-excitation block in
$\{\ket{N,g},\ket{N-1,e}\}$ is
\begin{equation}
 H_N=\begin{pmatrix}
 N\omega&g\sqrt N\\g\sqrt N&(N-1)\omega+\Omega
 \end{pmatrix},\quad
 E_{N,\pm}=N\omega+\frac{\Delta}{2}
       \pm\frac12\sqrt{\Delta^2+4g^2N}.
\end{equation}
The branches connected to the bare states are
\begin{align}
 E_g(n)&=n\omega+\frac\Delta2
    -\frac{\operatorname{sgn}\Delta}{2}\sqrt{\Delta^2+4g^2n},\\
 E_e(n)&=(n+1)\omega+\frac\Delta2
    +\frac{\operatorname{sgn}\Delta}{2}\sqrt{\Delta^2+4g^2(n+1)} .
\end{align}
The first formula also gives $E_g(0)=0$.
From $0\leq1+u/2-\sqrt{1+u}\leq u^2/8$ for $u\geq0$,
\begin{equation}
 |E_{\rm exact}-E_{\rm second\ order}|
 \leq g^4N^2/|\Delta|^3 .
\end{equation}
For $\omega=1,g=.03,\Delta=\pm.5$, exact $N=1,2,3$ blocks verify
this bound and both rate directions: the second-order ratios are
$1.00360649$ and $0.99640647$. These do not realize $\kappa=1/4$.

Beyond second order the clock gaps depend on $n$; the spectrum is not
an exact rigid rescaling. Bare $N_D$ is not conserved by $H_{\rm JC}$:
the effective conserved sectors are dressed, with correspondingly
transformed density and readout operators. Thus the microscopic check
supports the possible sign and dispersive structure, not an exact
identification of the bare two-level density with a universal clock rate.

\newpage
\section{Closure, physical interpretation and the next useful decision}
The requested sequence is complete for one declared candidate:
\begin{enumerate}
 \item A positive, ground-subtracted density operator and its reference
 measure are specified without assuming a spatial volume.
 \item One interaction is given. A common stationary constraint and
 matched finite state yield normalized joint clock probabilities.
 \item Relative rate follows from those probabilities under fixed
 calibration; common rescaling cannot create a differential effect.
\end{enumerate}
What has been obtained is an \emph{environment-dependent clock-frequency
shift}. Calling it universal slowing of time would require additional
physics. In particular, the coupling's sign and dependence on the clock
transition are not fixed by the density definition. Source fluctuations
can reduce visibility rather than define one deterministic time map.

\subsection*{What remains an assumption, not a hidden theorem}
Ordinary quantum mechanics, a subsystem split, the operational cell,
the phase POVMs, reference shielding, preparation and commensurate
spectra are inputs. No energy-information identity, negative energy,
mass conversion, emergent space, metric or universal proper time follows.
The finite reference cannot supply an exact canonical time operator.
Conditional phase arcs or record/epoch labels remain necessary at
recurrences, as established by G-R2.

The G-R2 finite-history recording gates and the present energy constraint
are \emph{not yet one autonomous energy-complete recording model}.
No detector design or physical memory lifetime is claimed. The exact
stationary test does not silently extend to arbitrary reference
interactions, irrational spectra, or an actual multilevel transmon.

\subsection*{Two physically useful follow-ups, not further fits}
First, select a source and clock realization that fixes the interaction
and its sign independently of a desired slowdown. If an intensive density
rather than total excitation energy is essential, compare differently
sized sources at fixed density with the same specified coupling profile.
Second, test whether independent clock transitions acquire the same
fractional shift. Transition-dependent distortions support ordinary
clock perturbation; a universal time interpretation needs a principle
explaining why these distortions disappear across different clocks.
Neither task is replaced by fitting a free function of $E_d$.

\subsection*{Executable verification}
The companion \texttt{verify\_gr3\_density\_clocks.py} passes 273 checks:
stationarity, positivity, conditional states, both POVMs, invariant
effects, two-clock likelihoods, relative slopes, energy-unit and
zero-shift invariance, the failed mean-density substitution, and
exact versus dispersive exchange spectra. Previous files are hashed.
All 4636 Route-G regression checks pass. These are algebra/implementation
checks, not experimental confirmation.

\begin{thebibliography}{9}
\raggedright
\bibitem{smith}
A.~R.~H.~Smith and M.~Ahmadi, \emph{Quantizing time: Interacting clocks
and systems}, Quantum \textbf{3}, 160 (2019),
\href{https://arxiv.org/abs/1712.00081v3}{arXiv:1712.00081v3}.
\bibitem{blais}
A.~Blais, A.~L.~Grimsmo, S.~M.~Girvin and A.~Wallraff,
\emph{Circuit quantum electrodynamics}, Rev.\ Mod.\ Phys.\
\textbf{93}, 025005 (2021), Sec.~III.B and Appendix~B.2.b,
\href{https://doi.org/10.1103/RevModPhys.93.025005}{doi:10.1103/RevModPhys.93.025005}.
\end{thebibliography}
\end{document}
